\subsection{半单与绝对半单自同态}

定义 6. --- 设 E 是交换域 K 上的有限维向量空间。则 E 的自同态 u 称为半单的，如果 E 的每个在 u 下封闭的子空间都有一个在 u 下封闭的补空间。

这意味着 K{[}X{]}-模 \(E_{u}\) 的每个子模都是直因子，换言之，K{[}X{]}-模 \(E_{u}\) 是半单的（VII，第 9 页）。

命题14. --- 设 u 是交换域上有限维向量空间的一个自同态，则 u 是半单的当且仅当 u 的最小多项式没有重因子。

这直接由 VII，第 9 页，推论 4 及第 31 页，注记 1 得出。

设 E 是交换域 K 上的向量空间，L 是 K 的一个扩域，u 是 E 的一个自同态；设 \(u_{(L)}\) 表示由 E 经标量扩张诱导的 L-自同态 \(1_{L} \otimes u\)（L-向量空间 \(\mathrm{E}_{(L)} = \mathrm{L} \otimes_{K} \mathrm{E}\) 上的）。同样地，若 F 是 E 的一个自同态集合，令 \(\mathfrak{F}_{(L)}\) 表示由 F 中的 u 所对应的 \(u_{(L)}\) 组成的集合。

推论。--- 设 u 是交换域 K 上有限维向量空间的一个自同态，并设 L 是 K 的一个扩域。若 \(u_{(L)}\) 是半单的，则 u 是半单的。若 u 是半单的且 L 在 K 上可分，则 \(u_{(L)}\) 是半单的。

这直接由命题 14 及 V，第 120 页，推论 1 得出（注意 u 与 \(u_{(L)}\) 的最小多项式重合）。

命题15. --- 设 E 是交换域 K 上的有限维向量空间，设 u 是 E 的一个自同态，并设 \(q(X)\) 为其最小多项式。则以下条件等价：

\begin{enumerate}
\def\labelenumi{(\roman{enumi})}
\item
  对每个扩张 \(\mathbf{L}\)（\(\mathbf{K}\) 的），自同态 \(u_{(\mathbf{L})}\) 是半单的。
\item
  存在一个扩张 \(\mathbf{L}\)（\(\mathbf{K}\) 的），使得自同态 \(u_{(\mathbf{L})}\) 是可对角化的。
\item
  多项式 \(q(\mathbf{X})\) 在 \(\mathbf{K}\) 上是可分的。
\end{enumerate}

事实上，条件 (i) 意味着：对 K 的任意扩域 L，多项式 \(1 \otimes q(\mathbf{X})\)（L{[}X{]} 中的）都没有重因子（命题 14）；条件 (ii) 意味着存在 K 的一个扩域 L，使得 \(q(\mathbf{X})\) 的所有根都属于 L 且这些根都是单根（VII，第 40 页，命题 12）；而由定义（V，第 38 页），这两个条件都等价于 (iii)。

定义7. --- 满足命题 15 的条件 (i)、(ii) 和 (iii) 的自同态 u 称为绝对半单的。

推论。--- u 是绝对半单的一个必要充分条件是：存在 K 的一个扩域 L，使得 L 是完全域且 \(u_{(L)}\) 是半单的。

推论中的条件意味着存在 K 的一个扩域 L，使得 L 是完全域且 \(q(\mathbf{X})\) 在 L{[}X{]} 中没有重因子（命题 14）；由 V，第 38 页，推论 2，此条件等价于 (iii)。

命题16. --- 设 E 是交换域 K 上的有限维向量空间，设 F 是 E 的一个自同态集合，并设 A 是 \(\mathrm{End}_{\mathbf{K}}(\mathbf{E})\) 中由 F 和 \(\mathrm{Id}_{\mathbf{E}}\) 生成的子代数。则以下条件等价：

\begin{enumerate}
\def\labelenumi{(\roman{enumi})}
\item
  存在一个扩张 \(\mathbf{L}\)（\(\mathbf{K}\) 的），使得 \(\mathfrak{F}_{(\mathbb{L})}\) 是可对角化的。
\item
  K-代数 A 是平展的（V，第 28 页，定义 1）。
\item
  \(\mathfrak{F}\) 的元素是绝对半单的，并且两两交换。
\end{enumerate}

首先注意，对 K 的每个扩张 L，由 \(\mathfrak{F}_{(L)}\) 与 \(Id_{E_{(L)}}\) 生成的 L-代数与 \(L \otimes_{K} A\) 重合；因此由命题 13，\(\mathfrak{F}_{(L)}\) 可对角化当且仅当 L-代数 \(L \otimes_{K} A\) 可对角化。于是条件 (i) 与 (ii) 的等价由 V，第 28 页，定义 1 得出。另一方面，(i) \(\Rightarrow\) (iii) 是显然的。最后，假设 (iii) 成立，并设 L 是 K 的一个代数闭包；则 \(\mathfrak{F}_{(L)}\) 的元素是可对角化的（VII，第 40 页，命题 12）且两两交换；因此由 VII，第 41 页，命题 13，\(\mathfrak{F}_{(L)}\) 是可对角化的。

推论。--- 两个可交换的绝对半单自同态的和与乘积都是绝对半单的。

注记。--- 设命题 16 的条件成立，并设 L 是 K 的一个扩域。由命题 13，集合 \(\mathfrak{F}_{(L)}\) 可对角化当且仅当代数 \(L \otimes_K A\) 可对角化。由 V，第 30 页，命题 2 可知，存在 K 的一个有限扩张 L 使 \(\mathfrak{F}_{(L)}\) 是可对角化的。事实上，L 还可以取为伽罗瓦扩张；的确，取一个有限子集 \(\mathfrak{F}'\)（\(\mathfrak{F}\) 的、生成 A 的），我们可以取 L 为 \(\mathfrak{F}'\) 中元素的最小多项式的一个分裂域（命题 12 和 13）。

